\documentclass[11pt]{article}
\usepackage[margin=1in]{geometry}
\usepackage[T1]{fontenc}
\usepackage{lmodern}
\usepackage{amsmath,booktabs,hyperref,microtype}
\hypersetup{hidelinks}
\makeatletter\let\tableinput\@@input\makeatother
% Generated; do not edit. Interval-robust, within-model.
\newcommand{\ClinAttackVersion}{17.1}
\newcommand{\ClinBackupUnionPct}{37.9}
\newcommand{\ClinCoreUnionPct}{5.3}
\newcommand{\ClinEmptyUnionPct}{44.9}
\newcommand{\ClinIdentityUnionPct}{4.2}
\newcommand{\ClinMinGuardFive}{3}
\newcommand{\ClinMinGuardNarrow}{3}
\newcommand{\ClinMinGuardOne}{7}
\newcommand{\ClinMinGuardTen}{2}
\newcommand{\ClinMinGuardWide}{14}
\newcommand{\ClinNamedGuardFive}{1}
\newcommand{\ClinServices}{4}


\title{Certified clinical-impact bounds for security-control selection:\\
a distribution-free method on an ATT\&CK-driven hospital model}
\author{}
\date{2026-09-20 (working note, version 0.1.0)}

\begin{document}
\maketitle

\begin{abstract}
Security-control decisions in hospitals are usually justified by abstract scores
(CVSS, coverage counts, reachability) that say nothing about patients. We pose and
solve a different problem: given a set of candidate defensive controls, certify ---
distribution-free, and robustly over acknowledged parameter uncertainty --- that
the probability of a sustained \emph{clinical-service} outage stays below a chosen
level. The model is a better synthetic hospital: its attack surface is the real
MITRE ATT\&CK Enterprise kill chain (v\ClinAttackVersion) with controls drawn from
ATT\&CK's own mitigations, and its damage is patient-facing outage of the
\ClinServices{} core clinical services through an explicit service-dependency
structure. Because every quantity is coordinatewise monotone in the model
parameters, adequacy is certified by evaluating the interval corners --- the
necessary/possible argument, lifted from efficiency to clinical outcome. The
undefended union-bound clinical outage is \ClinEmptyUnionPct\%, matching the
independently reported 44\% care-disruption rate. Identity controls dominate: a
four-control identity bundle certifies clinical outage $\le 5\%$, while backup,
segmentation, or patching alone each leave the outage near the undefended level.
The contribution is the method and problem formulation; effectiveness and
degradation are evidence-anchored intervals, not measurements of any hospital.
\end{abstract}

\section{Problem and model}
Let a control portfolio be a set of ATT\&CK mitigations. The adversary traverses
the real ATT\&CK kill chain; a portfolio lowers technique success through ATT\&CK's
\texttt{mitigates} edges. \emph{Impact reachability} $R$ is the probability of
completing the chain to the clinical-impact tactic (T1486 encryption, T1490 inhibit
recovery, T1489 service stop, T1485 destruction), computed as a product of
usage-weighted stage successes. A clinical service $s$ then suffers a sustained
outage with probability $R\cdot d_s$, where $d_s$ is a service-specific degradation
probability. Writing $O_s=R\,d_s$, we bound clinical impact two ways: per service,
and by the union bound $\sum_s O_s\ge \Pr(\ge 1\text{ service out})$.

Every parameter is an interval anchored to a published figure (base progression
from Mandiant/Sophos dwell data; mitigation effectiveness from the Microsoft MFA
study and a wide prior elsewhere; degradation $d_s\in[0.17,0.44]$ from Neprash 2022
and McGlave 2026). $O_s$ and the union bound are increasing in base and degradation
and decreasing in effectiveness.

\section{Certificate}
Monotonicity places the extremes of every bounded quantity at the corners of the
interval uncertainty set. A portfolio is \emph{guaranteed} clinically adequate at
level $\varepsilon$ when, at the adverse corner (base high, effectiveness low,
degradation high), every per-service outage \emph{and} the union bound are
$\le\varepsilon$; it is only \emph{possible} when this holds merely at the
favourable corner. The guarantee is distribution-free: it holds for every
parameter value in the set, with no distributional assumption. Where incident
counts exist, finite-sample betting bounds supply the intervals.

\section{Results}
Undefended clinical outage is \ClinEmptyUnionPct\% (union bound), a figure that
coincides with the independently observed 44\% care-disruption rate --- a coherence
check on the anchoring, not a fitted result. Along the certified cost-versus-outage
frontier, \ClinMinGuardTen{} controls guarantee union outage $\le 10\%$,
\ClinMinGuardFive{} guarantee $\le 5\%$, and \ClinMinGuardOne{} guarantee $\le 1\%$
(Table~\ref{tab:cfront}). Of the named bundles, only the identity bundle is
guaranteed at $5\%$ (union outage \ClinIdentityUnionPct\%); the hospital-core bundle
reaches \ClinCoreUnionPct\%, and single-family controls barely move the outcome ---
data backup alone leaves \ClinBackupUnionPct\% (Table~\ref{tab:cnamed}). The policy
reading is blunt and, to our knowledge, not something prior control-scoring methods
can state with a guarantee: for bounding \emph{patient-facing} outage, identity
controls dominate and no single other family suffices.

\begin{table}[t]\centering
\caption{Certified cost-versus-outage frontier (first eight greedy additions):
control added and worst-case union clinical outage after it, percent.}\label{tab:cfront}
\begin{tabular}{@{}rllr@{}}\toprule
Size & Mitigation & Name & Worst union outage (\%)\\\midrule
\tableinput table_clinical_frontier_rows.tex
\bottomrule\end{tabular}
\end{table}

\begin{table}[t]\centering
\caption{Named control bundles: worst-case union clinical outage and certificates.}\label{tab:cnamed}
\begin{tabular}{@{}lrrll@{}}\toprule
Bundle & Size & Worst union (\%) & Guar.\ $\le$10\% & Guar.\ $\le$5\%\\\midrule
\tableinput table_clinical_named_rows.tex
\bottomrule\end{tabular}
\end{table}

\section{Sensitivity to assumed uncertainty}
To guarantee union outage $\le 5\%$, a narrow effectiveness prior needs
\ClinMinGuardNarrow{} controls and a wide prior needs \ClinMinGuardWide{}: the
distribution-free guarantee demands more coverage precisely when less is known,
which is the honest behaviour a hospital decision-maker should want.

\section{Contribution and limitations}
The novelty is the formulation and method: a distribution-free certificate on
\emph{patient-facing clinical outage} for control portfolios, composing a real
ATT\&CK attack graph with a hospital service-dependency model, rather than scoring
abstract reachability. Limitations, stated plainly: ATT\&CK \texttt{mitigates} edges
are qualitative, not effect sizes; the impact-technique-to-service mapping, the
degradation probabilities, and the kill-chain progression are modelling choices
anchored to published ranges, not measured on a hospital; the model is synthetic
and unvalidated against incidents; and the union bound is deliberately conservative.
This is decision-support with an explicit guarantee under stated uncertainty, not an
empirical estimate of any hospital's risk.

\paragraph{Reproducibility.} MITRE ATT\&CK v\ClinAttackVersion{} is used unmodified
(Terms of Use), pinned and hashed; all outputs carry a provenance manifest and every
number here is generated from the committed tables.

\end{document}
